\section{Introduction}

Kecamatan Baureno, Kabupaten Bojonegoro, Jawa Timur, terletak di bantaran Sungai Bengawan Solo dan menghadapi ancaman banjir luapan rutin setiap musim penghujan. Wilayah ini mencakup 23 desa yang rentan terhadap genangan air akibat peningkatan debit sungai dari hulu dan curah hujan lokal yang deras. Sistem pelayanan kebencanaan yang efektif diperlukan untuk mendukung kesiapsiagaan warga dan koordinasi respons darurat, sejak tahap peringatan dini hingga evakuasi.

Permasalahan di lapangan mencakup tiga kendala utama: pertama, informasi banjir sering datang mendadak tanpa estimasi probabilitas risiko sebelum air meluap ke permukiman; kedua, warga mengalami kepanikan dan ketidaktahuan mengenai rute evakuasi aman saat banjir parah; ketiga, warga yang terjebak kesulitan meminta pertolongan kepada BPBD karena panggilan darurat jarang disertai koordinat lokasi akurat. Kendala tersebut menghambat respons cepat dan koordinasi penyelamatan yang efisien.

Untuk mengatasi masalah tersebut, dikembangkan SiagaBencana, sebuah platform web berbasis informasi hidrometeorologi yang mengintegrasikan data debit Sungai Bengawan Solo dan curah hujan untuk menghitung probabilitas banjir real-time (0--100\%) per desa. Sistem ini menyediakan peringatan dini notifikasi push, peta navigasi evakuasi otomatis yang hanya aktif saat banjir parah, serta kanal pelaporan darurat (SOS) dengan penangkapan GPS otomatis untuk mempercepat evakuasi warga terjebak. Platform ini dirancang untuk menghubungkan warga Baureno dengan petugas BPBD dalam ekosistem respons kebencanaan yang terkoordinasi.
